% Диаграмма Ганта — этапы курсовой работы
% Требует: \usepackage{pgfgantt} в преамбуле
\begin{figure}[htbp]
\centering
\resizebox{\textwidth}{!}{
\begin{ganttchart}[
    x unit=0.28cm,
    y unit title=0.6cm,
    y unit chart=0.7cm,
    vgrid={*6{draw=black!8}, *1{draw=black!25}},
    hgrid,
    title height=1,
    title label font=\footnotesize\bfseries,
    bar label font=\footnotesize,
    bar label node/.append style={left=0cm, anchor=east},
    bar height=0.5,
    bar/.append style={
        pattern=north east lines,
        pattern color=black!50,
        draw=black!60
    },
    time slot format=isodate,
    time slot unit=day
]{2026-04-09}{2026-05-28}
\gantttitle[title label font=\footnotesize\bfseries]{Апрель}{22}
\gantttitle[title label font=\footnotesize\bfseries]{Май}{28} \\
\gantttitle[title label font=\footnotesize]{9--15}{7}
\gantttitle[title label font=\footnotesize]{16--22}{7}
\gantttitle[title label font=\footnotesize]{23--30}{8}
\gantttitle[title label font=\footnotesize]{1--7}{7}
\gantttitle[title label font=\footnotesize]{8--14}{7}
\gantttitle[title label font=\footnotesize]{15--21}{7}
\gantttitle[title label font=\footnotesize]{22--28}{7} \\
\ganttbar{Этап 1: ТЗ}{2026-04-09}{2026-04-17} \\
\ganttbar{Этап 2: Диаграмма типов}{2026-04-17}{2026-04-24} \\
\ganttbar{Этап 3: Реализация}{2026-04-24}{2026-05-15} \\
\ganttbar{Этап 4: Отчёт и защита}{2026-05-15}{2026-05-22}
\ganttlink{elem0}{elem1}
\ganttlink{elem1}{elem2}
\ganttlink{elem2}{elem3}
\end{ganttchart}
}
\caption{Диаграмма Ганта~--- этапы выполнения курсовой работы}
\label{fig:gantt}
\end{figure}
